\begin{abstract}
Recursive training on model-generated data degrades a language model, and retaining
human-origin data in the mixture substantially reduces that degradation. This leaves a
question open: under a \emph{fixed lifetime budget} of human-origin optimizer tokens, when
during recursive training should those tokens be spent, and which under-covered modes of a
human reference distribution should they target? We predefine four budget-matched treatment
families --- random rescue, schedule-only, selection-only, and a joint time-and-mode policy
--- with complete recursive chains as experimental units, and freeze the allocation rule,
seed set, outcomes, and interpretation language before execution. The pipeline is validated
against a published positive control, reproducing its arm ordering and degradation
direction across two independent executions.

We execute \PilotChains{} chains over \PilotSeeds{} preregistered seeds and \PilotHorizon{}
generations under a corpus-assembly rule in which rescued human examples displace synthetic
records rather than being appended to them, so that the preregistered fairness constraint
--- equal lifetime human-origin tokens \emph{and} equal total optimizer tokens --- holds by
construction: human spend varies by \PilotHumanSpreadPct\% across arms and total optimizer
tokens are identical. Under those matched budgets the preregistered primary contrast is a
null. Joint time-and-mode allocation does not improve on the strongest non-joint baseline
(\PilotPrimaryMean{}, $95\%$ CI $[\PilotPrimaryLow, \PilotPrimaryHigh]$), and the interval
lies entirely inside the frozen practical equivalence region rather than merely spanning
zero. The secondary contrasts locate why: spending human tokens at all is worth
\PilotRandNonePct\% against a control trained on identical data volume, targeting
under-covered modes is worth a further \PilotSelRandPct\%, and \emph{when} within the chain
the budget is spent changes the primary outcome by \PilotSchedRandPct\% --- an interval
containing zero. At this operating point the allocation question decomposes: which modes the
budget targets matters, timing does not, and combining the two is equivalent to targeting
alone. Between-chain variance is small enough that the frozen seed set exceeds what the
preregistered threshold requires. Three of seven predeclared comparators, including a
non-deployable oracle upper bound, were not implemented, so how much headroom remains above
any of these policies is unknown.
\end{abstract}
